% 中文摘要
\chapter*{摘要}
\addcontentsline{toc}{chapter}{摘要}
递归自改进（Recursive Self-Improvement, RSI）是一类让模型通过
\emph{自我生成数据、自我评估、自我迭代}形成闭环训练的范式。在大
语言模型领域，自我精炼、自教学推理等机制已展示出无需额外人工标
注即可持续提升的能力；将这一范式引入序列推荐系统，有望缓解交互
数据稀疏、长尾曝光不足与标注获取昂贵等结构性瓶颈。然而，RSI 闭
环的已有实践暴露出两类系统性缺陷：其一，\textbf{自生成数据分布
漂移}——模型倾向于自我强化其头部预测，随迭代轮次累积放大，使训
练分布偏离真实交互分布；其二，\textbf{迭代误差累积}——缺少对自
生成样本质量的量化控制，低置信样本持续污染训练集，导致多轮迭代
后的性能崩塌。

本文以序列推荐为载体，系统研究 RSI 闭环中自生成数据的保真度控
制问题，主要工作与创新点如下：

\begin{enumerate}
\item \textbf{问题形式化与失败模式复现。}将 RSI 闭环建模为
"生成器—保真校验器—样本筛选—后继模型训练—评估反馈"的迭代过程，
给出分布漂移与误差累积的形式化刻画；在 MovieLens-100K 上以
朴素自训练（接受全部自生成样本）复现了迭代崩塌现象
（NDCG@10 由 \ExpNaiveStart{} 迭代 \ExpNaiveRounds{} 轮后跌至
\ExpNaiveEnd{}），证实崩塌源于头部自强化而非随机噪声。
\item \textbf{保真度门控机制。}提出一个双轴保真度门控器：
置信轴对候选样本的预测裕度按物品流行度分层做 z 标准化，消除置
信度对头部物品的结构性偏袒；分布轴以流行度分层上的
Jensen--Shannon 距离度量合成分布对真实分布的漂移，并按自适应
预算 $\varepsilon$ 在线收紧准入。门控仅放行同时满足置信与分
布预算的样本，从机制上切断头部自强化的正反馈。
\item \textbf{多视图一致性伪标签生成。}借鉴半监督学习中的
FixMatch 思想，为每个用户构造 $K$ 个扰动视图（物品随机丢弃与
尾部截断），仅当模型在各视图上的最高预测\emph{一致}且置信度、
裕度双阈值达标时，才生成伪标签样本。与朴素 top-$k$ 自生成相比，
一致性伪标签编码的是模型对扰动鲁棒的信念而非其头部偏好，为闭
环增益提供了信息来源。
\item \textbf{实验验证与可复现基础设施。}在 MovieLens-100K 的
留一法（leave-one-out）全量排序评测下，以 GRU4Rec 序列模型为骨
干、3 个随机种子开展系统实验，得到三个层次的结论：
（i）朴素自训练发生\emph{灾难性崩塌}——NDCG@10 由 \ExpNaiveStart{}
经 \ExpNaiveRounds{} 轮跌至 \ExpNaiveEnd{}（$p=\ExpNaivePBase{}$）
，同步的合成分布 JS 距离由 \NaiveJsFirst{} 扩大至 \NaiveJsLast{}
；（ii）本文框架将崩塌转化为\emph{受控衰减}——最终 NDCG@10 为
\ExpOursEnd{}，相对朴素自训练高 \ExpGainText{}（配对 $t$ 检验
$p=\ExpOursP{}$），混合分布漂移被预算锁定在 \OursJsRange{} 量级
（比朴素方案低约两个数量级）；消融实验呈单调阶梯——朴素
（\NaiveNdcg{}）$<$ 固定准入（\FixNdcg{}）$<$ 仅置信门控
（\ConfNdcg{}）$<$ 完整框架（\OursNdcg{}）——每个模块独立贡献
保护。（iii）\emph{诚实的边界}：所有 RSI 变体均未超过继续真实
训练的基线平台（\BaseNdcg{}，$p=\ExpOursPBase{}$），印证了自生
成数据的信息上界分析——在骨干模型已收敛的场景，RSI 的价值在
于可控性与安全性而非净增益。全部实验以种子固定、单线程确定性
的方式实现，并以指令流编排，结果可由命令行完整复现。
\end{enumerate}

\medskip
\noindent\textbf{关键词：}递归自改进；序列推荐；自生成数据；分布
漂移；保真度控制；半监督学习；伪标签；分布漂移预算

